%% 第6巻　未来を計算する　本文（第14回編集委員会）
%% 章は \bkinputchap{ch6-NN} で読む（まだない章は、組むときに飛ばして「BKPAGE missing」を出す）
%% 第6巻は巻の色が明るい青緑なので、部扉の絵の線は白でなく bkOn（濃い紺）でかく
%% Ver2（No.58）：章の前に Lua のごみ集め（第9巻と同じ。組版のメモリ対策）
\providecommand\bkinputchap[1]{\directlua{collectgarbage("collect") texio.write_nl("log","BKMEM #1 "..math.floor(collectgarbage("count")/1024).." MB")}\IfFileExists{#1.tex}{\input{#1}}{\clearpage\typeout{BKPAGE missing #1}}}
\bkinputchap{ch6-00}

%% ======================== 第1部 微分方程式の入口 ========================
\bkPartOpener{1}{第6巻\hspace{0.8em}第1部}{微分方程式の入口}%
  {172/{「次の瞬間」のルールだけで、未来がわかる？},
   160/{熱いお茶は、なぜだんだんゆっくり冷める？},
   148/{ばねにつけたおもりは、なぜ同じリズムでゆれる？}}%
  {% 少しずつ進む点（オイラー法）と、なめらかな曲線（濃い紺の線）
   \begin{scope}[bkOn, line width=1.2pt, line cap=round, line join=round]
     \draw (118,112) -- (172,112);
     \draw (122,108) -- (122,140);
     \draw[line width=1.6pt, domain=0:46, samples=50] plot ({122+\x}, {112+24*exp(-\x/14)});
     \foreach \k in {0,1,2,3,4,5} {\fill ({122+8*\k},{112+24*(1-8/14)^\k}) circle (1.5);}
   \end{scope}}%
  {「いまの状態から、次の瞬間どう変わるか」というルールを式にしたものが、微分方程式だよ。表計算で少しずつ進めて未来を計算し、冷めるお湯・電気回路・ばねの動きを読みとくよ。}
\bkPartContents{1}{%% 自動生成（gen.py）
\bktocchap{6-1}{微分方程式ってなに？}{答えが数ではなく「関数」になる方程式}{1}{2}{2}{26}{C}
\bktocchap{6-2}{表計算で解こう：オイラー法}{80 ℃のお湯は、10分後に何℃？}{1}{2}{2}{66}{C}
\bktocchap{6-3}{積分するだけで解ける形}{落ちるボールの位置を、式で求める}{2}{3}{3}{102}{C}
\bktocchap{6-4}{自分に比例してへる：冷めるお湯}{熱いほど速く冷める。答えは指数関数}{3}{3}{3}{132}{C}
\bktocchap{6-5}{空気抵抗と終端速度}{雨つぶは、なぜどこまでも速くならない？}{1}{2}{3}{176}{C}
\bktocchap{6-6}{ばねと振り子}{小さなゆれなら、周期はゆれの大きさによらない}{4}{3}{3}{206}{C}
\bktocchap{6-7}{ゆれながら止まる：減衰振動}{針がすっと止まるメーターのひみつ}{4}{3}{4}{246}{C}
}%
{物の動き・温度・電気の変化を微分方程式で書き、表計算で解いたり、積分で式の形の答えを出したりできるようになります。ばね・熱・電気の回路が同じ形の式になることも見えてきます。}
\bkpartbegin{1}{微分方程式の入口}
\bkinputchap{ch6-01}
\bkinputchap{ch6-02}
\bkinputchap{ch6-03}
\bkinputchap{ch6-04}
\bkinputchap{ch6-05}
\bkinputchap{ch6-06}
\bkinputchap{ch6-07}

%% ======================== 第2部 複素数と双曲線関数 ========================
\bkPartOpener{2}{第6巻\hspace{0.8em}第2部}{複素数と双曲線関数}%
  {172/{2乗してマイナスになる数は、何の役に立つ？},
   160/{かけ算で、点がくるりと回る？},
   148/{電線は、なぜこんな形にたるむ？}}%
  {% 複素平面の単位円と、回る矢印（濃い紺の線）
   \begin{scope}[bkOn, line width=1.2pt, line cap=round, line join=round]
     \draw (116,124) -- (172,124);
     \draw (144,106) -- (144,142);
     \draw (144,124) circle (14);
     \draw[-{Stealth[length=2.8mm]}, line width=1.6pt] (144,124) -- ({144+14*cos(40)},{124+14*sin(40)});
     \draw[line width=0.9pt] ({144+6},{124}) arc (0:40:6);
   \end{scope}}%
  {2乗して $-1$ になる数 $j$ を使うと、回転や波が、かけ算とたし算で計算できるようになるよ。オイラーの公式 $e^{j\theta}=\cos\theta+j\sin\theta$ が、指数関数と三角関数をつなぐよ。}
\bkPartContents{2}{%% 自動生成（gen.py）
\bktocchap{6-8}{2乗してマイナスになる数：虚数}{「2回かけて−1」になる数を、回転で考える}{4}{4}{4}{280}{C}
\bktocchap{6-9}{複素平面：かけ算は回転と拡大}{かけ算すると、矢印がくるっと回る}{4}{4}{4}{312}{C}
\bktocchap{6-10}{オイラーの公式}{指数関数と三角関数が、1本の式でつながる}{4}{4}{5}{344}{C}
\bktocchap{6-11}{双曲線関数 cosh・sinh}{たるませた電線の形は、放物線ではない}{3}{4}{5}{386}{C}
}%
{複素数を平面の点として計算し、オイラーの公式で回転と波を書けるようになります。電気の交流の計算（第9巻・物理シリーズ）と、ゆれの式を解く道具になります。}
\bkpartbegin{2}{複素数と双曲線関数}
\bkinputchap{ch6-08}
\bkinputchap{ch6-09}
\bkinputchap{ch6-10}
\bkinputchap{ch6-11}

%% ======================== 第3部 式で解く ========================
\bkPartOpener{3}{第6巻\hspace{0.8em}第3部}{式で解く}%
  {172/{車のサスペンションは、なぜゆれをすぐに止める？},
   160/{ブランコをこぐと、なぜ大きくゆれる？},
   148/{物理の法則は、みんな同じ形で書ける？}}%
  {% 減衰振動の曲線と、その包絡線（濃い紺の線）
   \begin{scope}[bkOn, line width=1.2pt, line cap=round, line join=round]
     \draw (118,122) -- (172,122);
     \draw[line width=1.6pt, domain=0:50, samples=120] plot ({120+\x}, {122+14*exp(-\x/18)*cos(\x*25)});
     \draw[line width=0.8pt, dash pattern=on 2pt off 2pt, domain=0:50, samples=30] plot ({120+\x}, {122+14*exp(-\x/18)});
     \draw[line width=0.8pt, dash pattern=on 2pt off 2pt, domain=0:50, samples=30] plot ({120+\x}, {122-14*exp(-\x/18)});
   \end{scope}}%
  {指数関数 $e^{st}$ を入れてみる、という1つのアイデアで、ゆれながら止まる動きや共振が式で解けるよ。いくつもの量がからみあう動きも、和と差に分けると解けるよ。}
\bkPartContents{3}{%% 自動生成（gen.py）
\bktocchap{6-12}{減衰振動を式で解く：特性方程式}{2次方程式の判別式が、ゆれ方を決める}{4}{5}{5}{420}{C}
\bktocchap{6-13}{ゆらし続けると：強制振動と共振}{小さな力でも、タイミングが合うと大きくゆれる}{4}{5}{5}{458}{C}
\bktocchap{6-14}{微分方程式がいくつも：連立微分方程式}{つながった2つの水そうの水位は、どう変わる？}{5}{5}{5}{492}{C}
\bktocchap{6-15}{まとめ：物理は「微分方程式＋はじめの状態」で書ける}{ばねも電気も熱も、同じ形の式で書ける}{1}{2}{3}{534}{C}
}%
{特性方程式で減衰振動を解き、共振のしくみを説明できるようになります。第9巻のラプラス変換と制御は、この部の考え方をそのまま使います。}
\bkpartbegin{3}{式で解く}
\bkinputchap{ch6-12}
\bkinputchap{ch6-13}
\bkinputchap{ch6-14}
\bkinputchap{ch6-15}

%% ======================== 巻末 ========================
\bkinputchap{ch6-back}
